\subsection{A Basic Income, and Which Floor Leaves a Record}
\label{sec:8.3.4}
The competing floor is a basic income, and the choice between them is usually argued on forecasts about how much work machines will take. It does not have to be. Both proposals are conditional, differently, and the difference is what matters: a basic income is conditional on continued political willingness to fund it, which the aggregation machinery can withdraw in a budget without producing a record of anything; a jobs guarantee is conditional on a placement decision, which is a lever, and a dangerous one under a hostile administration, but a lever that leaves a decision-maker, a record, a counterparty and an appeal behind it.

That is section~\ref{sec:9.3.5}'s slope instrument — measure whether the cost of dissent is rising, not whether a channel exists — turned on this book's own proposal. Dependence that generates artifacts is safer than dependence that generates none, not because the first is benign but because the first can be measured and the second cannot. It is also, for the same reason, the more demanding proposal to build: everything in the comparison rests on the appeal being real, and the previous section says what makes it real.

That standard now has a court behind it, and on the half of the case that travels. When the Supreme Court kept a Federal Reserve governor in her post in 2026 against a removal for cause, what it held she was owed was not the post but the artifact: “some explanation of the evidence at issue, some avenue for a response, and a deadline by which a response would be due” \autocite{cook2026}. What does not travel is this: the Fed's protection rests on a pedigree no body chartered later can acquire. What a removal has to produce before it counts is not tied to that pedigree.

I put a jobs guarantee ahead of a basic income on that ground. One remaining measure does work the floor cannot: broadened social insurance — unemployment benefits, healthcare, housing — reaches the people who cannot work or choose not to, whom a guarantee by construction does not reach at all.
